\subsection{A product form of the total}
\label{sec:comp-total-energy-product}
\begingroup
\def\comptonproductlabel{eq:comp-total-energy-product}
Products of elementary logarithms can make the collected result shorter without introducing additional polylogarithms. Define $L=\ln(1+X)$ and $M=\ln X$. The two GPLs with the same coefficient combine into
\begin{equation}
\mathcal P(X)=G(0,-1,0;X)+G(0,1,0;X)=\tfrac12\left[\operatorname{Re}\operatorname{Li}_3(X^2)-M\operatorname{Re}\operatorname{Li}_2(X^2)\right].
\end{equation}
For the derivation, first take $0<X<1$. The shuffle relation $G(0)G(0,a)=2G(0,0,a)+G(0,a,0)$, with $G(0,a;X)=-\operatorname{Li}_2(X/a)$ and $G(0,0,a;X)=-\operatorname{Li}_3(X/a)$, gives
\begin{equation}
 G(0,a,0;X)=2\operatorname{Li}_3(X/a)-M\operatorname{Li}_2(X/a),\qquad a=\pm1.
\end{equation}
Summing the two signs and using $\operatorname{Li}_n(z)+\operatorname{Li}_n(-z)=2^{1-n}\operatorname{Li}_n(z^2)$ gives the stated result. Both shuffle multiplicities and this duplication factor are needed for its factor $1/2$. For $X>1$, continue all functions on the same edge of the principal cut. On either edge, the imaginary parts of $\operatorname{Li}_2(X)$ and $\operatorname{Li}_3(X)$ are respectively $\pm\pi\ln X$ and $\pm\pi\ln^2X/2$, with the same sign. These follow by integrating their derivatives from $X=1$, and they cancel in $2\operatorname{Li}_3(X)-M\operatorname{Li}_2(X)$. The analogous cancellation holds for the paired expression at argument $X^2$. Thus both edges give the same real function and the continuation adds no branch constant. The following equivalent form uses only real polylogarithms with arguments in $(0,1]$:
\begin{equation}
\mathcal P(X)=\begin{cases}\tfrac12[\operatorname{Li}_3(X^2)-M\operatorname{Li}_2(X^2)],&0<X\leq1,\\ \tfrac12[\operatorname{Li}_3(X^{-2})+M\operatorname{Li}_2(X^{-2})]+\tfrac13M^3+\tfrac16\pi^2M,&X\geq1.\end{cases}
\end{equation}
To derive the second line, put $t=\ln X>0$, $D(t)=\operatorname{Li}_2(e^{-2t})$ and $T(t)=\operatorname{Li}_3(e^{-2t})$. Their derivatives are $D^{\prime}(t)=2\ln(1-e^{-2t})$ and $T^{\prime}(t)=-2D(t)$. The identity $\ln(e^{2t}-1)=2t+\ln(1-e^{-2t})$ then gives
\begin{align}
\operatorname{Re}\operatorname{Li}_2(e^{2t})&=\frac{\pi^2}{3}-2t^2-D(t),\\
\operatorname{Re}\operatorname{Li}_3(e^{2t})&=T(t)-\frac43t^3+\frac{2\pi^2}{3}t.
\end{align}
Indeed, the derivatives of the two sides are respectively $-2\ln(e^{2t}-1)$ and $2\operatorname{Re}\operatorname{Li}_2(e^{2t})$. Write $\zeta(n)=\sum_{k=1}^{\infty}k^{-n}$ for the Riemann zeta function at integers $n>1$. Their values at $t=0$ agree because $\operatorname{Li}_2(1)=\pi^2/6$ and $\operatorname{Li}_3(1)=\zeta(3)$. These derivative and boundary conditions fix both identities and give $\mathcal P(1)=\zeta(3)/2$.
The sole apparent rational pole at $X=1$ can also be removed before evaluation. Define the continuous elementary function
\begin{equation}
R_{\log}(X)=\begin{cases}\ln X/(X-1),&X\ne1,\\1,&X=1,\end{cases}\qquad R_{\log}(1+h)=1-\frac h2+\frac{h^2}{3}+O(h^3).
\end{equation}
The series follows by dividing the Taylor series of $\ln(1+h)$ by $h$, which converges for $|h|<1$. Extracting the common power $X^{-3}$ gives the exact total coefficient as
\begin{equation}
C_1(X)=\frac{1}{X^3}\left[R_{\log}(X)+\mathcal E(X,L,M)+\sum_{j=1}^{13}b_j(X,Y,L,M)\mathcal F_j(X)\right].
\label{\comptonproductlabel}\end{equation}
The elementary part is
\begin{equation}
\mathcal E=e_{0}L^{3}+e_{1}L^{2}M+e_{2}L^{2}+e_{3}LM+e_{4}L+e_{5}M+e_{6}.
\end{equation}
\begin{equation}
e_{0}=- \frac{X^{2} - 16 X - 23}{3}.
\end{equation}
\begin{equation}
e_{1}=4 X^{2} - 2 X - 12.
\end{equation}
\begin{equation}
e_{2}=- 2 X^{3} - \frac{11 X^{2}}{2} + \frac{133 X}{2} + 36 + \frac{96}{X + 4} - \frac{6}{X + 1} + \frac{1}{\left(X + 1\right)^{2}}.
\end{equation}
\begin{equation}
e_{3}=X^{2} + 5 X - 10 + \frac{12}{X + 1} - \frac{4}{\left(X + 1\right)^{2}}.
\end{equation}
\begin{equation}
e_{4}=\frac{37 X^{2}}{4} - 32 X - 30 - \frac{15}{X + 1} - \frac{1}{\left(X + 1\right)^{2}}.
\end{equation}
\begin{equation}
e_{5}=- \frac{14 X^{5} + 31 X^{4} + 14 X^{3} - 6 X^{2} - 8 X - 2}{2 \left(X + 1\right)^{3}}.
\end{equation}
\begin{equation}
e_{6}=- \frac{X \left(273 X^{3} - 982 X^{2} - 2960 X - 1744\right)}{24 \left(X + 1\right)^{2}}.
\end{equation}
The thirteen remaining functions and their coefficients are
\begin{equation}
\mathcal F_{1}=G(0,-1;X).
\end{equation}
\begin{equation}
b_{1}=X^{3} + X^{2} - 43 X + 6 - \frac{96}{X + 4} - \frac{6}{X + 1} + \frac{2}{\left(X + 1\right)^{2}}- 2 L X \left(2 X + 5\right).
\end{equation}
\begin{equation}
\mathcal F_{2}=G(-2,-1;X).
\end{equation}
\begin{equation}
b_{2}=- 4 X^{2} - 80 X - 68 - \frac{128}{X + 4} + \frac{8}{X + 1} - \frac{4}{\left(X + 1\right)^{2}}.
\end{equation}
\begin{equation}
\mathcal F_{3}=G(1,0;X),\qquad b_{3}=- X^{2} + X - 2.
\end{equation}
\begin{equation}
\mathcal F_{4}=G(r,-1;X).
\end{equation}
\begin{equation}
b_{4}=Y \left(3 X^{3} + 18 X^{2} + 44 X - 8 - \frac{64}{X}\right)- 3 M Y \left(X - 2\right) \left(X + 4\right).
\end{equation}
\begin{equation}
\mathcal F_{5}=G(0,-1,-1;X),\qquad b_{5}=- 5 X^{2} - 8 X + 32.
\end{equation}
\begin{equation}
\mathcal F_{6}=G(-1,-2,-1;X),\qquad b_{6}=8 \left(X^{2} - 4 X - 6\right).
\end{equation}
\begin{equation}
\mathcal F_{7}=G(-4,r,-1;X),\qquad b_{7}=Y \left(- 6 X^{2} + 72 + \frac{48}{X}\right).
\end{equation}
\begin{equation}
\mathcal F_{8}=\mathcal P(X),\qquad b_{8}=2 X^{2} + 4 X - 12.
\end{equation}
\begin{equation}
\mathcal F_{9}=G(r,0,-1;X),\qquad b_{9}=Y \left(- 5 X^{2} + 14 X + 88 + \frac{48}{X + 4} + \frac{48}{X}\right).
\end{equation}
\begin{equation}
\mathcal F_{10}=G(r,-1,0;X),\qquad b_{10}=- 3 Y \left(X - 2\right) \left(X + 4\right).
\end{equation}
\begin{equation}
\mathcal F_{11}=G(r,-1,-1;X),\qquad b_{11}=Y \left(15 X^{2} - 18 X - 216 - \frac{96}{X + 4} - \frac{96}{X}\right).
\end{equation}
\begin{equation}
\mathcal F_{12}=G(r,-2,-1;X),\qquad b_{12}=Y \left(- 16 X^{2} + 32 X + 256 + \frac{64}{X + 4} + \frac{192}{X}\right).
\end{equation}
\begin{equation}
\mathcal F_{13}=G(r,r,-1;X),\qquad b_{13}=9 X^{2} - 24.
\end{equation}

\endgroup
This form follows from the complete coefficient table using only the identities
\begin{align}
G(-1,0)&=LM-G(0,-1),\\
G(-1,0,-1)&=L\,G(0,-1)-2G(0,-1,-1),\\
G(-1,-1,0)&=\tfrac12ML^2-L\,G(0,-1)+G(0,-1,-1),\\
G(0,r,-1)&=M\,G(r,-1)-G(r,0,-1)-G(r,-1,0).
\end{align}
The common endpoint is $X$. In addition, repeated letters $-1$ give powers of $L$ divided by factorials. Expanding every product back by the shuffle identity reproduces all twenty-three original coefficients exactly. The two words combined into $\mathcal P$ have the same coefficient $c_{15}=c_{16}$, so this last grouping replaces their sum without another factor of two.
The extraction of $R_{\log}$ is fixed by the coefficient $c_2$ of $M$ in the collected table. Direct polynomial division gives
\begin{equation}
 X^3c_2(X,Y)=\frac{1}{X-1}+e_5(X),\qquad X^3c_2(X,Y)M=R_{\log}(X)+e_5(X)M.
\end{equation}
The second equality at $X=1$ means its continuous limit. All other coefficients are multiplied by $X^3$ before the common factor $X^{-3}$ is placed outside the bracket. Thus this extraction changes neither the normalization of $C_1$ nor the coefficient of its removable term. Every displayed rational coefficient is regular at $X=1$, as follows by exact polynomial cancellation. The function $G(1,0;X)=\zeta(2)-\operatorname{Li}_2(1-X)$ is also explicitly real and regular for every $X>0$. Its derivative is $\ln X/(X-1)$ and its limit as $X\to0^+$ is zero, which fixes the identity without numerical continuation.
